% 复制本文件用于后续分问时，修改题号、所有 q1 标签和对应引用。
\section{问题一的模型建立与求解}
\label{sec:q1-model}

\subsection{数据与建模依据}
\slot{说明输入文件、字段、单位、处理规则和处理前后样本量。由数据特征或机理关系解释模型选择；参数说明来源，缺失信息明确标记。}

\subsection{模型建立}
\slot{先说明现实关系，再给出数学表达，随后解释符号、适用域和用途。写清变量、目标、约束、参数来源及边界条件，使公式与代码口径一致。}
% 公式示例见 examples/（默认不加载）；此处只填写本题实际模型。
% 使用 equation/align 环境，并设置 \label{eq:q1-model} 等唯一标签。

\subsection{求解方法与实现}
\slot{说明算法、关键参数、初始条件、停止准则、随机种子、运行环境和脚本入口。区分可行解、启发式最好解与有最优性证明的解，准确报告求解状态。}

\subsection{结果与分析}
\slot{先直接回答本问，再用图表支持答案并解释原因；比较必须具有共同条件和指标。表中数字、单位和口径对应真实结果台账。}
表\ref{tab:q1-results}预留本问核心答案，尚未计算的项目保持待填。
\begin{table}[htbp]
  \centering
  \caption{问题一的核心结果（待填）}
  \label{tab:q1-results}
  \small
  \begin{tabularx}{\linewidth}{lYY}
    \toprule
    输出 & 结果及单位 & 条件与解释 \\
    \midrule
    \slot{输出项} & \slot{真实结果} & \slot{条件与题意} \\
    \bottomrule
  \end{tabularx}
\end{table}

\subsection{模型验证}
\slot{根据本问主要风险选择可行性、量纲守恒、样本外误差、基线比较、扰动或随机误差检查。写明实际执行的方法、范围、结果与证据，不机械套用全部检验。}
